\section{Proyectos: reservas e itinerarios}

La aplicación \textbf{Proyectos} organiza el trabajo operativo tras una venta (o en paralelo): tareas por etapa, adjuntos de vouchers y responsables.

\subsection{Abrir Proyectos}

\begin{enumerate}
    \item Pulse \textbf{Proyectos} en el menú de aplicaciones.
    \item Seleccione el \textbf{proyecto} configurado para itinerarios o reservas (nombre definido en la implementación).
    \item La vista \textbf{Kanban} muestra columnas (etapas). Cada tarjeta es una \textbf{tarea} o \textbf{elemento de trabajo}.
\end{enumerate}

\subsection{Crear o abrir una tarea de reserva}

\begin{enumerate}
    \item Pulse \textbf{Nuevo} o el \textbf{+} en la columna adecuada del tablero.
    \item Asigne un \textbf{Título} claro (por ejemplo: ``Reserva MP -- Juan Pérez -- 12--18 junio'').
    \item En \textbf{Asignados}, elija el usuario responsable si aplica.
    \item En \textbf{Descripción} o \textbf{Notas}, detalle fechas, pasajeros y observaciones operativas.
    \item Pulse \textbf{Guardar} o \textbf{Crear} según el botón visible.
\end{enumerate}

\subsection{Adjuntar documentos en la tarea}

\begin{enumerate}
    \item Abra la tarea.
    \item En el \textbf{chatter} o zona de adjuntos, cargue \textbf{vouchers}, \textbf{PDF de boletos}, \textbf{confirmaciones de hotel}, etc.
    \item Mueva la tarjeta arrastrándola a la siguiente \textbf{etapa} cuando el trabajo avance (por ejemplo hacia compra de tickets o ejecución).
\end{enumerate}

\subsection{Enlace con el CRM y Ventas}

Mantenga coherencia: el \textbf{contacto} y los montos acordados deben coincidir con la \textbf{oportunidad ganada} y el \textbf{pedido} en Ventas. Si cambian fechas o pasajeros, actualice primero la ficha del cliente y los documentos de venta, luego la tarea del proyecto.
